
\chapter*{Preface}

This report presents the research and development carried out as part of our final year group project on the interpretable analysis of 3D chest CT images. The project addresses two limitations of current deep learning based approaches to automated abnormality detection: their reliance on post-hoc explainability techniques that provide only approximate, loosely traceable localizations, and the absence of a unified way to benchmark models across classification, segmentation, and explainability simultaneously. To this end, the project makes two contributions. First, we propose an inherently interpretable dictionary learning framework that learns compact volumetric atoms representing normal lung structures and disease-related radiological patterns, allowing detected abnormalities to be back-projected to the original voxel space and traced directly to the anatomical structures that drove the diagnostic decision. Second, we develop a unified evaluation framework that benchmarks state-of-the-art models on a publicly available 3D chest CT dataset across three axes: classification accuracy, segmentation quality, and explainability.

The report covers the theoretical background, the proposed dictionary learning method, its implementation, and the design of the evaluation framework, followed by an analysis of the experimental results. Our method is benchmarked against state-of-the-art deep learning baselines using this framework, and the results show that it achieves competitive classification and segmentation performance while offering markedly improved interpretability.

We would like to thank our supervisors, collaborators, and everyone who provided data, code, guidance, and feedback during this project. Any errors or omissions in this report are the responsibility of the project team.